% Glossary terms from ch17: lines of the form \item[Term] Definition.
\item[Near] The set $\Near(V,\state,r)=\set{\vect{v}\in V:\norm{\vect{v}-\state}\le r}$ of tree vertices within the connection radius of a point; the candidate set of \textsc{ChooseParent} and \textsc{Rewire} in \rrtstar.
\item[Connection radius] $r_n=\min\{\gamma_{\rrtstar}(\log n/n)^{1/d},\eta\}$ with $n$ the number of tree vertices; asymptotic optimality needs $\gamma_{\rrtstar}>2(1+1/d)^{1/d}(\mu(\Xfree)/\zeta_d)^{1/d}$, and the schedule keeps the expected \textsc{Near} set at $\bigO{\log n}$ vertices.
\item[ChooseParent] The \rrtstar step that connects the new vertex to the near vertex giving the smallest cost-to-come through a collision-free segment, with the nearest vertex as fallback.
\item[Rewire] The \rrtstar step that makes the new vertex the parent of every near vertex whose cost-to-come it lowers through a collision-free segment, propagating the saving to all descendants of that vertex.
\item[Cost-to-come] $\mathrm{Cost}(\vect{v})$, the length of the tree path from the start to vertex $\vect{v}$; \rrtstar keeps $\mathrm{Cost}(\vect{v})=\mathrm{Cost}(\mathrm{parent}(\vect{v}))+\norm{\vect{v}-\mathrm{parent}(\vect{v})}$ as an invariant.
\item[Asymptotic optimality] The property that the cost of the best path in the tree converges to the optimal cost $c^*$ with probability one as the number of samples grows; \rrtstar has it, \rrt does not.
\item[Robustly optimal solution] An optimal path that is the limit in cost of collision-free paths with positive clearance; the assumption under which sampling-based planners can converge to it.
\item[Strong $\delta$-clearance] A path all of whose points are at distance at least $\delta$ from the obstacles.
\item[$k$-nearest RRT*] The variant of \rrtstar whose \textsc{Near} set is the $k_n=\lceil k_{\mathrm{RRG}}\log n\rceil$ nearest vertices with $k_{\mathrm{RRG}}>e(1+1/d)$; it needs no estimate of $\mu(\Xfree)$.
\item[Informed set] $\set{\state:\norm{\state-\state_{\mathrm{init}}}+\norm{\state-\state_{\mathrm{goal}}}\le c_{\mathrm{best}}}$, the only region in which a path cheaper than the current best can pass; the analog of the region $\fcost\le C^*$ of \astar.
\item[Prolate hyperspheroid] The shape of the informed set: an ellipse (2D) or ellipsoid of revolution (3D) with foci at the start and the goal, transverse diameter $c_{\mathrm{best}}$ and conjugate diameter $\sqrt{c_{\mathrm{best}}^2-c_{\min}^2}$.
\item[Informed RRT*] \rrtstar that, once a solution of cost $c_{\mathrm{best}}$ exists, samples the informed set uniformly by mapping the unit ball through $\state=\mat{C}\mat{L}\state_{\mathrm{ball}}+\state_{\mathrm{center}}$.
\item[Branch-and-bound pruning] Deleting tree vertices with $\mathrm{Cost}(\vect{v})+\norm{\state_{\mathrm{goal}}-\vect{v}}\ge c_{\mathrm{best}}$, which can never lie on a better path; used in anytime \rrtstar.
\item[BIT*] Batch informed trees: samples are drawn in batches inside the informed set and the implicit random geometric graph is searched with an edge queue ordered by $\hat f=\hat g+\hat h$ with lazy collision checks, reusing the search across batches.
